\section{Introduction}
\label{sec:introduction}

Agentic AI systems are rapidly evolving from passive question-answering assistants into autonomous decision-makers that plan, call tools, and act in open-ended environments~\cite{mialon2023gaia,drouin2024workarena,xie2024osworld}. This shift is especially visible in information-retrieval (IR) contexts, where modern agents now \emph{act} on retrieved documents and structured data---issuing follow-up queries, writing files, sending messages, and modifying records---rather than merely returning passages~\cite{xi2024agentir}. As these systems become embedded in real-world workflows, their failures are no longer limited to incorrect answers on static tasks. Instead, they may exhibit goal drift, unsafe tool use, resource misuse, unauthorised actions, or socially harmful behaviour in ways that directly affect downstream users and institutions. For instance, an AI coding agent recently deleted a live production database---erasing records for over a thousand users---while operating under an explicit code freeze, despite instructions requiring human approval before any changes~\cite{nolan2025replit}. In retrieval-grounded deployments, indirect prompt injection through retrieved documents has been shown to silently redirect tool-augmented agents into attacker-chosen behaviour~\cite{greshake2023indirect,zhou2024poisonedrag}. Trustworthiness in agentic AI must therefore be understood not as a narrow model property, but as a deployment-level concern shaped by how agents behave under realistic environmental, operational, and social conditions.

Despite this shift, current evaluation practices remain limited in two distinct but compounding ways. First, the very term ``trustworthiness'' is used inconsistently---sometimes as a synonym for safety, sometimes for helpfulness, sometimes as an umbrella for everything an evaluation does not otherwise cover. Existing surveys list many desiderata~\cite{liu2024trustllm,wang2024trustworthyagent} but do not give a small, measurable set of properties against which an agentic system can be assessed. Without such a definition, ``HAAF passes a trustworthiness audit'' would itself be a vacuous claim. Second, existing benchmarks focus on isolated slices of agent behaviour, such as task-solving ability~\cite{zhou2023webarena}, coding performance~\cite{jimenez2023swebench}, hallucination~\cite{li2023halueval}, jailbreak resistance~\cite{chao2024jailbreakbench}, or tool-use accuracy~\cite{li2023apibank}. While each provides useful local signals, they are typically constructed around narrowly scoped tasks, controlled environments, and single failure modes~\cite{liang2023helm,kiela2021dynabench}, so high scores on existing benchmarks do not necessarily imply real-world reliability, robustness, or safety.

We argue that the central limitation is therefore not merely that current evaluations cover too few dimensions, but that they lack \emph{both} (i)~a concrete definition of trustworthiness and (ii)~a principled notion of \emph{representativeness}. Agent trustworthiness is not a point property that can be inferred from a handful of benchmark instances; rather, it is a distributional property that depends on how an agent behaves over a heterogeneous, heavy-tailed space of socio-technical scenarios. When the distribution of benchmark tasks is misaligned with the distribution of real deployment conditions, evaluation results can create a false sense of confidence~\cite{koh2021wilds,shirali2022dynamicbench}. This problem is particularly severe for low-frequency but high-consequence events, including adversarial manipulations, cascading tool failures, and socially sensitive interactions, which are often absent from standard evaluation suites despite being central to trustworthy deployment.

To address both gaps, we propose the \textbf{Holographic Agent Assessment Framework (HAAF)}.
Our key idea is to first \emph{define} agent trustworthiness as a five-property profile (\S\ref{subsec:defining_trustworthiness})---Reliability, Robustness, Safety, Social-Ethical Alignment, and Operational Integrity---and then to \emph{measure} that profile over a \emph{scenario manifold} that spans task types, tool interfaces, interaction dynamics, social contexts, and risk levels, rather than over disconnected benchmark islands. The framework centres on a \emph{distribution-aware representative sampling engine} that defines \emph{what} should be tested, while three complementary probing interfaces---static cognitive and policy analysis, interactive sandbox simulation, and social-ethical alignment assessment---define \emph{how} agents should be evaluated. These components are connected through an iterative \emph{Trustworthy Optimization Factory}: red-team probing exposes vulnerability surfaces, blue-team hardening designs targeted interventions, and re-evaluation verifies improvement, cycling until the deployed system meets deployment-readiness thresholds. Rather than evaluating model weights in isolation, we assess \emph{deployed agentic systems} comprising the model, prompting policies, tool interfaces, guardrails, and oversight mechanisms.

We illustrate this paradigm with four concrete pieces of evidence: a coverage audit of existing benchmarks against the five properties; a Factory cycle on a single focal agent that reduces violation rate from 16.7\% to 4.2\% with no new improper refusals; a \emph{cross-model trustworthiness profile} over 13 contemporary agentic systems spanning seven model families (Llama, Mistral, Kimi, GLM, Qwen, GPT, DeepSeek) served through both local and Amazon Bedrock back-ends; and a \emph{cross-family Factory generalisation experiment} in which the focal-model interventions are applied uniformly to those same 13 systems and shown to transfer rather than over-fit. The cross-model profile surfaces property-level trade-offs---one system stronger on robustness, another on social-ethical alignment---that any single aggregate metric would erase. In summary, this paper makes the following four contributions:
\begin{itemize}[leftmargin=*]
\item We give an operational five-property definition of agentic trustworthiness (P1--P5), grounded in current AI risk frameworks~\cite{nist2023airmf,eu2024aiact}, and use it as the explicit measurement target of HAAF.
\item We identify \emph{representativeness} as the central missing principle in current agent trustworthiness evaluation, arguing that the primary challenge lies in distribution design rather than solely in metric design.
\item We propose HAAF, a unified framework for representative, multi-view evaluation over risk-sensitive scenario distributions, and a comparison methodology that supports principled multi-property contrasts between agentic systems.
\item We outline an iterative evaluation-hardening methodology and demonstrate its feasibility through (a)~a 24-scenario Factory cycle on a single focal model, (b)~a 100-scenario cross-model trustworthiness profile over 13 contemporary agentic systems from seven families, and (c)~a cross-family transfer experiment in which the focal-model interventions reduce risk for all 13 systems without per-model or per-scenario tuning.
\end{itemize}
